\section{Simultaneous savings for balanced cuts}\label{sec:lamps}

The circuit decomposition in Lemma~\ref{lem:balanced} allows each summand
to use a different cut of the rows. To bound the whole sum, we need one
placement of the rows at which \emph{every} cut admits a saving. Each
side of a cut gets its own supplier. The two suppliers must be close
enough to share the rank index at low cost. Moving them away from the
cone vertex must shorten the legs to their rows in total.

Fix a nonempty class \(\mathcal M\) of Hadamard spaces without ports,
closed under positive rescaling, closed convex subspaces, pointed
ultralimits and finite \(\ell_2\) products, as supplied by
Theorem~\ref{prod:main}. Write \(X=[n]\), \(n\ge2\), and \(m=\log_2 n\).
All logarithms and entropies in this section use base two. A
\emph{balanced layout} is a map \(b:X\to\{0,1\}\) taking each value on
at least \(n/4\) rows. In particular, every balanced cut in
Lemma~\ref{lem:balanced} defines such a layout.

On a metric cone with vertex \(0\) set
\(
 \langle x,z\rangle=(\|x\|^2+\|z\|^2-d(x,z)^2)/2
\).
This pairing is homogeneous in each variable under nonnegative dilations.
It adds coordinatewise in products of cones.

\begin{theorem}[Binary lamp certificate]\label{lamp:certificate}
For every list of \(L\) balanced layouts on \([n]\) there is a Hadamard
metric cone in \(\mathcal M\), with unit points \(x_1,\ldots,x_n\)
of barycenter zero and two points \(Z_{\ell,0},Z_{\ell,1}\) per layout,
such that
\[
 d(Z_{\ell,0},Z_{\ell,1})\le1,\qquad
 \|Z_{\ell,v}\|\le C\sqrt{m+\log L},\qquad
 \sum_{i=1}^n\langle x_i,Z_{\ell,b_\ell(i)}\rangle
     \ge c n/\sqrt m
\]
for every layout. The constants are absolute.
\end{theorem}

For the cost comparison in Section~\ref{sec:comparison}, move the two
suppliers from the vertex to
\(\tau Z_{\ell,0},\tau Z_{\ell,1}\). The wire between them has length
at most \(\tau\). For a unit row site \(x_i\), the initial rate at
which its leg shortens is \(\langle x_i,Z_{\ell,b_\ell(i)}\rangle\).
The last estimate guarantees a total first-order saving for each layout,
and the radius bound controls the quadratic error.

To construct the row sites, first fix an arbitrary distribution on the
layouts and construct a pair of points for each layout. Row \(i\) selects one point from every
pair according to its layout bit. Write \(\operatorname{bar}\) for the
Cartan barycenter, the minimizer of the mean squared distance.

\begin{lemma}[Barycentric separation]\label{lamp:lemma}
Let \(b_1,\ldots,b_L\) be balanced layouts, with any distribution
on \([L]\). In a member of \(\mathcal M\) there are points
\(z_{\ell,0},z_{\ell,1}\) satisfying
\[
 d(z_{\ell,0},z_{\ell,1})\le1,\qquad
 \operatorname{diam}\{z_{\ell,v}\}\le C\sqrt{m+\log L}.
\]
For \(u_i=\operatorname{bar}_{\ell} z_{\ell,b_\ell(i)}\) and
independent uniform \(i,i'\in X\),
\(
 (\mathbb E d(u_i,u_{i'})^2)^{1/2}\ge c/\sqrt m
\).
The constants are absolute.
\end{lemma}

To prove Lemma~\ref{lamp:lemma}, we encode a word as the increment of
one lamp. A layout reports the sides of selected letters. By successively hashing
the unread information into shorter words, we arrange that the combined
reports identify a typical increment up to a short list, for any
prescribed sequence of layouts. The lamp graph charges less for an increment on
this list than for an arbitrary increment.

Applying Theorem~\ref{thm:entropy} to the lamp graph gives a placement in
which changing an unprepared lamp has large mean squared length. We compare this change
with a path that first prepares the reports and then makes the same
increment. The difference in lengths forces the row barycenters to
separate. The tangent cone at their barycenter turns separation into
the first-order saving, and a product of cones realizes the saving for
all layouts at once.

\subsection{Identifying a word from its reports}

A uniform letter of \(X\) contains \(m\) bits of information, and a
layout reads one binary digit from it. We use batches of letters. At
each stage a report conveys about half a bit per letter in the batch.
A fixed map sends the batch to a shorter one whose distribution, given
a typical report, is close to uniform.

Choose a small absolute \(\delta\in(0,1/100)\), whose size will be
fixed below, and put
\[
  r=\lceil K_\delta(m+\log L)\rceil,\qquad S=rm,\qquad r_1=r.
\]
Given \(r_a\), set
\(r_{a+1}=\lfloor(1-(1/2+\delta)/m)r_a\rfloor\), stopping before
the next length would be smaller than \(\delta r\). Denote the last
index by \(P\) and put \(R=\sum_{a=1}^P r_a\). Geometric decay gives
\[
  \delta r\le r_P\le3\delta r,\qquad P=O_\delta(m),\qquad R\le2S,
  \qquad \sum_a\sqrt{r_a}\le C m\sqrt r.
\]
We take \(K_\delta\) large enough for the rounding estimates, for
\(P(m+1)\le\delta S\), and for \(P\log L\le S\). Thus \(S\) is the
information in the uniform initial word, and \(R\) will be the storage
allocated to its successive reports.

Let \(H_\ell\) be the binary entropy of \(b_\ell(i)\) for uniform
\(i\). Balance gives \(H_\ell>1/2\). The report
\(F_{a,\ell}(y)\) of \(y\in X^{r_a}\) records the layout bits on
the first \(\lfloor r_a/(2H_\ell)\rfloor\) letters, padded by zeros
to length \(r_a\). A report of positive probability \(p\) under the
uniform word is typical if
\( |{-\log p}-r_a/2|\le\delta r_a/4\).
The information content is a sum of independent bounded variables with
mean within one of \(r_a/2\). Thus the probability of an atypical
report is at most \(\exp(-c_\delta r_a)\).

For each \(a<P\) choose a map
\(\sigma_a:X^{r_a}\to X^{r_{a+1}}\) as follows. Assign its values
independently and uniformly. Inside any typical report cell and for
each proposed output word, the expected fiber size lies between
\(2^{3\delta r_a/4}\) and \(2^{5\delta r_a/4+m}\), by the definition
of \(r_{a+1}\). The multiplicative Chernoff bound puts every fiber within
a factor \(1\pm2^{-\delta r_a/8}\) of its mean except with probability
at most \(2\exp(-c2^{\delta r_a/2})\) per test. A union bound covers
the at most \(L2^{O(mr_a)}\) choices of layout, report and output,
since \(r_a\ge\delta r\) and \(K_\delta\) is sufficiently large.
Fix maps satisfying these estimates for all the layouts at each step.

A \emph{chart} is a list \(\theta=(\ell_1,\ldots,\ell_P)\) of layouts,
one for each batch. For a uniform \(X_1\in X^r\), put
\(X_{a+1}=\sigma_a(X_a)\). Its report at stage \(a\) is
\(F_{a,\ell_a}(X_a)\). We use two properties, both uniform in the chart.
\begin{enumerate}
\item The joint law of \(X_a\) and all reports before time \(a\)
is within \(\exp(-c'_\delta r)\) in total variation of the product
of a fresh uniform batch and the reports of independent uniform
batches. Indeed, on typical cells the next hash is nearly uniform
by the fiber estimate. The probability of the other cells and the
relative fiber errors add at each step under composition of the
report and hash kernels. The resulting factor \(P\) is absorbed in
the exponential. In particular all reports are typical with the same
error bound after decreasing \(c'_\delta\).
\item If all reports are typical, at most \(2^{8\delta S}\) initial
words give these reports. Guess the last batch at cost at most
\(3\delta S\) bits and work backwards. Each fiber then costs at
most \(1+5\delta r_a/4+m\) bits. The claimed bound follows from
\(R\le2S\) and \(P(m+1)\le\delta S\).
\end{enumerate}

\subsection{The lamp graph and its placement}

Identify \(X^r\) with a cyclic group \(A\) of order \(2^S=n^r\)
by any bijection, so that each word is a possible increment. A state
\((g,\theta,h)\) consists of an array \(g\in A^M\) of lamp values, a
chart \(\theta\), and \(R\) report bits \(h=(h_1,\ldots,h_P)\). For
each chart and report, the associated list consists of the initial words
whose successive batches give exactly those reports. The graph has the
following reversible edges. The cost of an edge is a bound on its
\emph{squared} length.
\begin{itemize}
\item A single report bit can be flipped at cost \(\log(1/\delta)\).
A layout and its report in one batch can be reset at cost \(r\).
Resetting the entire chart and report costs \(S\).
\item One can add any word to a specified lamp at cost \(S\).
When all reports are typical, one can add or subtract a word from
their list at cost \(\delta S\).
\end{itemize}
The number \(M\) of lamps will be an absolute constant. Fix a reference
chart and report, and call the state with those data and array \(g\)
the reference state of \(g\). The demand law \(\nu\) joins the
reference states of independent uniform arrays. Even given the
starting state, its destination array has \(MS\) bits of entropy.

A physical walk can change the array only through the edges just
described. To compare this with the independent arrays of a demand
step, we put the two walks in the setting of Theorem~\ref{thm:entropy}.
For a symmetric edge law \(\alpha\) on a finite connected weighted
graph, let \(E_\alpha\) be its mean cost.
The \emph{padded physical walk} has stationary pair law equal to
one half \(\alpha\) plus one half the diagonal law of the demand marginal.
The padded demand walk is defined with the two roles interchanged.
Both are the stationary Markov chains of their respective pair laws
and have the same stationary distribution. For a pair law \(\mu\) write
\(\mathcal E_\mu(f)=\mathbb E_{(u,v)\sim\mu}d(f(u),f(v))^2\).
The entropy rate of a stationary array process \((G_t)\) is
\(\lim_{t\to\infty}H(G_1,\ldots,G_t)/t\).

\begin{lemma}[Placing the lamps]\label{lamp:entropy-placement}
Suppose that all edge costs lie in \([1,S]\), the lamp-array demand
is as above, and the entropy rate of the array along every padded
physical walk is at most \(C_0 E_\alpha\). For sufficiently large
\(M\), depending only on \(C_0\), the graph admits a placement in
\(\mathcal M\) with every squared edge length at most its cost and
mean squared demand distance at least \(c_0 MS\), where
\(c_0>0\) depends only on \(C_0\).
\end{lemma}
\begin{proof}
Let \(B\) be the maximum mean squared demand distance under these edge
constraints. The possible squared distance tables form a closed convex
cone by the product, scaling and ultralimit hypotheses. Connectedness
makes the constrained tables bounded. For any table with demand
energy one we have \(\max_e d_e^2/c_e\ge1/B\), where \(c_e\) is the
edge cost. Convex minimax against the probability simplex on the
edges gives weights \(w_e\) such that
\(\sum_e w_e d_e^2/c_e\ge1/B\) for every such table.
Weight edge \(e\) proportionally to \(w_e/c_e\) and orient it uniformly
to obtain \(\alpha\). Homogeneity now gives the squared energy comparison
\(\mathcal E_\nu\le(B/E_\alpha)\mathcal E_\alpha\) for all placements
in \(\mathcal M\). The nontriviality needed to normalize demand energy
follows since a class of one-point spaces has ports.

Padding divides both energies by two. In the padded demand walk the
conditional entropy of the destination array given the starting state is
at least \(MS/2\), as is seen by conditioning also on the coin selecting
the demand or the diagonal step. By Theorem~\ref{thm:entropy}, with comparison constant \(1+B/E_\alpha\),
\[
 MS/2\le C(1+B/E_\alpha)(1+C_0 E_\alpha)\le C'(S+B).
\]
The last inequality uses \(1\le E_\alpha\le S\). Choosing \(M\ge4C'\)
proves the claim.
\end{proof}

We verify the coding hypothesis with an absolute \(C_0\). Encode the
initial full state once in a long physical path and then the move type,
including a hold, with a bounded number of bits per step. On a reset
encode the new chart and report for the affected batches. For a single
batch this uses \(O(r)\) bits including its address. For a full reset it
uses \(O(S)\), since \(P\log L\le S\) and \(R\le2S\).
On single-bit flips transmit only the move type. Sending a bit address
at every flip would cost \(\lceil\log R\rceil\) bits per flip. Instead
the decoder keeps the report last received, and the encoder sends a
correction when a lamp uses it.

At the next lamp move first correct the decoder's copy. A discrepancy
of Hamming weight \(D\) costs at most
\(1-D\log\delta-(R-D)\log(1-\delta)
=O(1+\delta R+D\log(1/\delta))\) bits by a prefix code for independent
Bernoulli discrepancies. The sum of the values \(D\) over the path
is at most the number of single flips, since corrections refresh the copy
and resets overwrite the corresponding entries. The remaining data are
the lamp address and either the whole increment or its signed index
on the current list. Taking \(\delta S\ge1+\log M\), the list bound pays
for these data at each cheap move. The accumulated description length
is at most an absolute constant times the accumulated unpadded move costs
plus the length of the path and the initial description. If several moves realize an observed
transition, sample the move conditionally from the edge law defining that transition. Each such
description decodes to the correct array sequence, and its expected
physical cost is the one charged by \(E_\alpha\). Taking the stationary
expectation, dividing by path length and taking the limit bounds the
entropy rate by \(C_0E_\alpha\). In this estimate \(C_0\) is independent
of \(\delta\). Thus \(M\) may be fixed absolutely, and \(K_\delta\)
then made large enough for the address bound.

\subsection{Changing one lamp through the batches}

Let \(p(g,\theta,h)\) be the placement furnished by
Lemma~\ref{lamp:entropy-placement}. For a fixed sequence
\(y=(y_1,\ldots,y_P)\) of batches take independent layouts with the
given law, let \(h_a=F_{a,\ell_a}(y_a)\), and average the point
\(p(g,\theta,h)\) by successive Cartan barycenters in the order
\(\ell_P,\ell_{P-1},\ldots,\ell_1\). Call the result \(T(g,y)\).
Each average holds all earlier layouts fixed. This order allows us to
replace a batch after revealing only the preceding reports.
Barycenters contract mean distance under any coupling
\cite{Sturm}, hence also root-mean-square distance.

For the calculations below we recall the coordinate-exchange inequality.
If \((V_i,V'_i)\) are independent exchangeable pairs, and \(V^{(i)}\)
uses \(V'_i\) in coordinate \(i\) and \(V\) elsewhere, then
\begin{equation}\label{lamp:exchange}
 \mathbb E d(f(V),f(V'))^2
 \le\sum_i\mathbb E d(f(V),f(V^{(i)}))^2
\end{equation}
for any Hadamard-valued map \(f\), by Equation~\eqref{prod:exchange}.

We compare the same increment twice, once with independent batches
and once with the reports of the increment itself. In each comparison
we couple the layouts identically. The order of the barycenters then
allows us to apply Equation~\eqref{lamp:exchange} to individual letters.

The full-reset move puts \(T(g,y)\) within \(\sqrt S\) of its reference
state for every \(y\). For independent uniform \(g,g'\), their mean-square
separation after taking barycenters therefore has square root at least
\(\sqrt{c_0MS}-2\sqrt S\). Enlarge \(M\) absolutely if necessary.
Applying Equation~\eqref{lamp:exchange} to the \(M\) array entries shows that
for a uniform lamp index \(j\) and an independent uniform word \(X_1\)
the root-mean-square distance from \(T(g,y)\) to
\(T(g+e_jX_1,y)\) is at least \(c_1\sqrt S\). Here \(e_jX_1\) denotes
the increment at lamp \(j\). This holds also after choosing independent
uniform batches \(Y=(Y_a)_a\), independently of \(g,j,X_1\).

Write \(Y^{(a)}=(X_1,\ldots,X_a,Y_{a+1},\ldots,Y_P)\), so
\(Y^{(0)}=Y\), and set \(X^*=(X_1,\ldots,X_P)\). Figure~\ref{lamp:fig-path}
shows the path. Replace independent batches by those of \(X^*\), one
at a time. With the entire chain installed, its typical reports put
\(X_1\) on their list, so the squared cost of adding \(X_1\) to lamp
\(j\) is at most \(\delta S\). On atypical reports we can still use
the edge of cost \(S\). Averaging over the same charts bounds the
root-mean-square length of the increment by
\(\sqrt{\delta S+S\exp(-c'_\delta r)}\).

\begin{figure}[tbp]
\centering
\begin{tikzpicture}[font=\small,
  state/.style={draw=ink!65,fill=ink!4,rounded corners=2pt,
    minimum height=20pt,minimum width=76pt,inner sep=3pt},
  path step/.style={-{Latex[length=2mm]},draw=ink,semithick}]
\node[state] (a) at (0,0) {$T(g,Y)$};
\node[state] (b) at (8.0,0) {$T(g,X^*)$};
\node[state] (c) at (8.0,-1.85) {$T(\widetilde g,X^*)$};
\node[state] (d) at (0,-1.85) {$T(\widetilde g,Y)$};
\node[above=3pt of a] {Independent batches};
\node[above=3pt of b] {Batches of the increment};
\draw[path step] (a.east) -- node[above]{replace batches \(1,\ldots,P\)} (b.west);
\draw[path step,draw=sea] (b.south) -- node[left=3pt]{\(O(\sqrt{\delta S})\)} (c.north);
\draw[path step] (c.west) -- node[below]{restore batches \(P,\ldots,1\)} (d.east);
\draw[densely dashed,draw=ink] (a.south) -- node[right=3pt]{\(\ge c_1\sqrt S\)} (d.north);
\end{tikzpicture}
\caption{The same lamp increment \(\widetilde g=g+e_jX_1\) with two
choices of reports. The vertical labels give root-mean-square distances.
The dashed segment compares the endpoints of the three-step detour.
With the increment's own batches in place, the middle step uses the
short list with high probability.}
\label{lamp:fig-path}
\end{figure}

Let \(J\) be the supremum of the root-mean-square separation in
Lemma~\ref{lamp:lemma} under its stated distance bounds, with \(C\)
chosen sufficiently large in terms of \(K_\delta\).
We claim that the root-mean-square cost of replacing batch \(a\) is at most
\begin{equation}\label{lamp:path-estimate}
 \sqrt{r_a\log(1/\delta)}\,J+\exp(-c''_\delta r).
\end{equation}
First average over layouts \(\ell_P,\ldots,\ell_a\), holding the
earlier layouts fixed. The points for \(Y^{(a-1)}\) and \(Y^{(a)}\)
depend on \(X_1,\ldots,X_{a-1}\) only through their common prefix
reports. The unaltered \(Y_a,\ldots,Y_P\) are independent
of the entire hash chain. By the joint law above, in the
expectation of this squared distance we can replace \(X_a\) by a fresh
independent uniform batch and use independent prefix reports, with
error \(O(S\exp(-c'_\delta r))\), since all the points concerned lie in a
ball of radius \(\sqrt S\) at the same reference state. This estimate
averages over prefix reports and requires no estimate conditioned on
one specified report value.

Apply Equation~\eqref{lamp:exchange} to the two independent words at batch \(a\).
Fix the other letters and the prefix reports. Immediately before the
average over \(\ell_a\), the letter in question affects only the bit
read by \(b_{\ell_a}\). If this position is ignored by that layout, the
two choices give the same point. At each layout the two points are at
distance at most \(\sqrt{\log(1/\delta)}\), by the flip move and coupling
all subsequent layouts. Over all layouts these points have diameter
at most \(\sqrt r\), by a one-batch reset with that same coupling.
After rescaling by \(1/\sqrt{\log(1/\delta)}\), their averages over
\(\ell_a\) are precisely a test for \(J\). Each summand in
Equation~\eqref{lamp:exchange} is thus at most \(J^2\log(1/\delta)\).
Completing the outer barycenters contracts the mean square. Taking
square roots and absorbing polynomial factors into the exponential
gives Equation~\eqref{lamp:path-estimate}. On the return path the translated
array is again uniform independently of all batches, since \(g\) was
independent uniform.

Sum the replacement lengths along both horizontal paths in
Figure~\ref{lamp:fig-path}. The \(L^2\) triangle inequality gives
\[
 c_1\sqrt S\le
 2J\sqrt{\log(1/\delta)}\sum_a\sqrt{r_a}
 +\sqrt{\delta S+S\exp(-c'_\delta r)}
 +2P\exp(-c''_\delta r).
\]
Choose \(\delta\) sufficiently small compared with \(c_1^2\), and then
\(K_\delta\) sufficiently large. Since \(\sum_a\sqrt{r_a}\le C m\sqrt r\)
and \(S=rm\), this gives \(J\ge c/\sqrt m\). Choosing a placement
within a factor two of the supremum proves Lemma~\ref{lamp:lemma}.

\subsection{From barycenters to directions}

\begin{proof}[Proof of Theorem~\ref{lamp:certificate}]
We first normalize the mean squared row norm to one. Given a distribution
on layouts, take the points of Lemma~\ref{lamp:lemma}, put
\(p=\operatorname{bar}_i u_i\), and set
\(\sigma^2=n^{-1}\sum_i d(p,u_i)^2\). The triangle inequality gives
\(\sigma\ge c/(2\sqrt m)\). In the completed tangent cone at \(p\) take
\[
 x_i=\log_p(u_i)/\sigma,\qquad Z_{\ell,v}=\log_p(z_{\ell,v}).
\]
To realize the tangent cone in \(\mathcal M\), rescale distances
by integers \(k\to\infty\). Represent a direction of length \(t\)
by the points at distance \(t/k\) on its geodesic. In a pointed
ultralimit this is an isometric embedding and extends to the
completed cone. Its image is closed and convex. The logarithm map preserves distance to \(p\)
as norm and contracts pairwise distances. First variation of the mean
squared distance at its minimizer shows that
\(\sum_i\langle v,\log_p u_i\rangle\le0\) for each represented direction,
and hence for every tangent vector by continuity. Thus the \(x_i\)
have barycenter zero \cite{BH,Sturm}. The bounds on the \(Z\)'s follow
from contraction and from the diameter bound on the convex hull of
the original points, which contains \(p\).

Use the variance inequality for the measure with barycenter \(u_i\),
evaluated at \(p\). It gives
\[
 \mathbb E_\ell d(p,z_{\ell,b_\ell(i)})^2
 \ge d(p,u_i)^2+\mathbb E_\ell d(u_i,z_{\ell,b_\ell(i)})^2 .
\]
On expanding the cone pairing and using contraction, we obtain
\(\mathbb E_\ell\langle x_i,Z_{\ell,b_\ell(i)}\rangle
\ge d(p,u_i)^2/\sigma\). The total expected score is therefore at
least \(n\sigma\). Write \(s_\ell=\sum_i\langle x_i,Z_{\ell,b_\ell(i)}\rangle\) for
the score at layout \(\ell\). For every distribution \(q\) there is
now a configuration with \(\sum_\ell q_\ell s_\ell\ge cn/\sqrt m\).
For each configuration the strict inequality
\(\sum_\ell q_\ell s_\ell>cn/(2\sqrt m)\) defines an open set of
distributions. These sets cover the compact simplex. A finite subcover
gives a finite matrix, with configurations as rows, layouts as columns,
and the scores as entries. Minimax gives a convex combination of its
rows with score at least \(cn/(2\sqrt m)\) in every column. Realize this combination by
an \(\ell_2\) product, multiplying the vectors in each factor by the
square root of its weight. All metric bounds and the normalization
survive, and barycenters in the product are coordinatewise.

Finally enlarge the original list to include every cyclic translate
\(i\mapsto b_\ell(i+t)\) of every layout, with indices modulo \(n\).
Apply the preceding construction to the enlarged list, indexing each
layout by \((\ell,t)\). In the product of \(n\) copies take the rows
\(n^{-1/2}(x_{i+t})_t\) and centers
\(n^{-1/2}(Z_{(\ell,-t),v})_t\). At row \(i+t\) layout
\((\ell,-t)\) reads the original bit \(b_\ell(i)\). Thus the row
at position \(i\) has squared norm \(n^{-1}\sum_k\|x_k\|^2=1\).
The product rows are centered and each layout keeps its score bound. The enlarged list has at most \(nL\) members,
so \(m+\log(nL)\le2(m+\log L)\) preserves the radius estimate.
\end{proof}
